% Official source: https://www.imo-official.org/assets/documents/problems/2011/2011_eng.pdf
\competitionyear{2011}

\begin{problem}{2011}{1}
For a set $A=\{a_1, a_2, a_3, a_4\}$ of four distinct positive integers, put $s_A=a_1+a_2+a_3+a_4$, and let $n_A$ be the number of pairs $(i, j)$ with $1\le i<j\le4$ for which $a_i+a_j$ divides $s_A$. Determine all such sets $A$ for which $n_A$ is as large as possible.
\end{problem}

\begin{problem}{2011}{2}
Let $S$ be a finite set of at least two points in the plane, with no three collinear. A \emph{windmill} begins with a line $\ell$ passing through exactly one point $P\in S$. Rotate $\ell$ clockwise about $P$ until it first passes through another point $Q\in S$. The point $Q$ becomes the pivot, and the line continues rotating clockwise about $Q$ until it reaches another point of $S$. Repeat this process indefinitely.

Prove that $P$ and $\ell$ can be chosen so that every point of $S$ becomes a pivot infinitely many times.
\end{problem}

\begin{problem}{2011}{3}
Suppose that $f\colon\R\to\R$ satisfies
\[
    f(x+y)\le yf(x)+f(f(x))
\]
for all $x, y\in\R$. Prove that $f(x)=0$ whenever $x\le0$.
\end{problem}

\begin{problem}{2011}{4}
A balance and $n$ weights of masses $2^0, 2^1, \dots, 2^{n-1}$ are given, where $n$ is a positive integer. Place the weights on the balance one at a time, choosing at each step an unused weight and putting it on either pan. The right pan must never be heavier than the left pan, including after each placement.

How many ways are there to place all $n$ weights under this restriction?
\end{problem}

\begin{problem}{2011}{5}
Let $f\colon\Z\to\N$ be a function such that $f(m-n)$ divides $f(m)-f(n)$ for every $m, n\in\Z$. Prove that whenever $m, n\in\Z$ satisfy $f(m)\le f(n)$, we have $f(m)\mid f(n)$.
\end{problem}

\begin{problem}{2011}{6}
Let $ABC$ be an acute triangle with circumcircle $\Gamma$. Take a tangent line $\ell$ to $\Gamma$, and reflect $\ell$ in the lines $BC$, $CA$, and $AB$ to obtain $\ell_a$, $\ell_b$, and $\ell_c$, respectively. Prove that the circumcircle of the triangle formed by $\ell_a$, $\ell_b$, and $\ell_c$ is tangent to $\Gamma$.
\end{problem}
